(1) 近轴光线以高度和倾角列向量 \((x,\theta)^T\) 表示。传播距离 \(d\) 的矩阵为 \(P(d)=\begin{pmatrix}1&d\\0&1\end{pmatrix}\)，焦距 \(f\) 的薄透镜矩阵为 \(L(f)=\begin{pmatrix}1&0\\-1/f&1\end{pmatrix}\)。从物面出发，光线依次经历 \(P(f)\)、第一透镜、\(P(2f)\)、第二透镜和 \(P(f)\)，矩阵乘积须按光路从右向左排列。
(2) 总矩阵 \(\mathsf M=P(f)L(f)P(2f)L(f)P(f)=-\mathsf I\)，因此像面高度 \(x_{\rm out}=-x_{\rm in}\)，出射倾角 \(\theta_{\rm out}=-\theta_{\rm in}\)。成像矩阵的 \(B=0\)，所以输出高度不依赖输入光线倾角，物面上任意一点被成像为唯一位置；横向放大率 \(A=-1\)，恰是等倍倒像。
(3) 第一透镜后焦面位于两个透镜中间，入射物场的不同近轴平面波分量由第一镜会聚到不同位置。若横向波数为 \(k_x\)、总波数为 \(k_0=2\pi/\lambda\)，该平面坐标 \(X\simeq f k_x/k_0=f\lambda k_x/(2\pi)\)。所以在此平面直接裁剪的不是物体位置，而是其角谱或空间频率，理想薄透镜还附带可预测的二次相位因子。
(4) 傅里叶面开口 \(|X|<a\) 等价于保留 \(|k_x|<k_{\rm cut}=2\pi a/(f\lambda)\) 的频谱。尖锐边缘和细小纹理对应高空间频率，会被抑制，像面变平滑且衍射点扩散函数展宽，典型可分辨尺度为 \(O(1/k_{\rm cut})\)。一个振幅光阑会同时改变复频谱的模与可观察相位关系，因此不能宣称仅仅降低亮度而不改变成像结构。
